% CLAUDE REPORT::WCDC-0001::2026-11-15T04:51-08:00
% Worst-case minimum derivation cost: first audit of the shared directive.
\documentclass[11pt]{article}
\usepackage{amsmath,amssymb,amsthm}
\usepackage[margin=1in]{geometry}
\usepackage{booktabs}
\usepackage{hyperref}
\newtheorem{theorem}{Theorem}
\newtheorem{lemma}[theorem]{Lemma}
\newtheorem{proposition}[theorem]{Proposition}
\newtheorem{corollary}[theorem]{Corollary}
\theoremstyle{definition}
\newtheorem{definition}[theorem]{Definition}
\newcommand{\Wb}{\overline W}
\newcommand{\poly}{\operatorname{poly}}
\newcommand{\cost}{\operatorname{cost}}
\title{WCDC-0001: Worst-Case Minimum Derivation Cost $\Wb_R(N)$ and the\\ Weakest Sufficient Bridge $\mathrm{SAT}\in\mathrm{P}\Rightarrow\Wb_R\le\poly$}
\author{Claude scheduled routine (worst-case derivation directive)}
\date{2026-11-15 04:51 America/Los\_Angeles}
\begin{document}
\maketitle

\begin{abstract}
This run starts the shared directive
\texttt{research/coordination/2026-11-15\_0451\_WORST\_CASE\_DERIVATION\_COST\_SHARED\_DIRECTIVE.md}.
The target is the worst-case minimum derivation cost $\Wb_R(N)$ of a fixed inference
system $R$ on a multisign DAG hypergraph encoding of 3-CNF. We prove four exact
classification statements about the proposed bridge
$\mathrm{SAT}\in\mathrm{P}\Rightarrow \Wb_R(N)\le\poly(N)$, together with one computational
hostile control. First, $\Wb_R$ is the Cook--Reckhow proof-length function of $R$. Second,
for every $R$ with a known superpolynomial lower bound, the bridge is \emph{equivalent}
to $\mathrm{P}\ne\mathrm{NP}$, so it is a reformulation and not a weaker step. Third, for
every $R$ that p-simulates Extended Frege, the bridge is equivalent to a precise
statement: every correct polynomial-time SAT decider has polynomial-size $R$-proofs of
its own correctness tautologies. We call this the \emph{correctness-provability arrow}.
It is the exact open arrow. Fourth, a pointwise transfer ``every algorithm run is a cheap
$R$-derivation'' fails for resolution-type reductions on Tseitin formulas. We check this
by exact minimum-width computations against GF(2) elimination counts. The
proof spine has two OPEN load-bearing arrows. \textbf{P versus NP remains OPEN.}
\end{abstract}

\section{Objects and quantifiers}
\begin{definition}
A \emph{derivation system} $R$ over the encoding $\langle F\rangle$ is a polynomial-time
checker $V_R(\langle F\rangle,\pi)$ that is sound (accepts only when $F$ is
unsatisfiable) and complete (every unsatisfiable $F$ has some accepted $\pi$). The cost
$\cost_R(\pi)\ge|\pi|$ charges every hyperedge, sign, shared DAG node and extension
symbol. Set
\[
\Wb_R(N)=\max_{F\in\mathrm{UNSAT},\ |\langle F\rangle|\le N}\ \min_{\pi:\,F\vdash_R\bot}\cost_R(\pi).
\]
\end{definition}
The multisign DAG hypergraph and the ``minimal reduced logical state'' fit this
definition. Each reduction step rewrites signed hyperedges. A reduction sequence ending
in the contradictory state $\bot$ is a derivation $\pi$, and its checker is local. The
minimum over sequences is exactly the nondeterministic choice in Cook--Reckhow.

\paragraph{Quantifiers for the separation route.}
The route would be: (LB) $\forall k\ \exists^\infty N:\ \Wb_R(N)>N^k$; (BR)
$\mathrm{SAT}\in\mathrm{P}\Rightarrow\exists k\ \forall N:\Wb_R(N)\le N^k+k$. Then
(LB)$\wedge$(BR)$\Rightarrow\mathrm{P}\ne\mathrm{NP}$. This uses one fixed $R$ in both.

\section{Exact classification results}
\begin{theorem}[KNOWN THEOREM, Cook--Reckhow 1979]\label{thm:cr}
Some derivation system $R$ has $\Wb_R(N)\le\poly(N)$ if and only if
$\mathrm{NP}=\mathrm{coNP}$.
\end{theorem}
\noindent\emph{Consequence.} Granting the optimal sequence for free is legitimate.
It does mean the measured quantity is \emph{nondeterministic} refutation length. A
superpolynomial lower bound for every $R$ is $\mathrm{NP}\ne\mathrm{coNP}$, which is
stronger than $\mathrm{P}\ne\mathrm{NP}$ as far as is known. So the directive fixes one
$R$, which is the right choice.

\begin{theorem}[PROVED/EXACT, elementary]\label{thm:reform}
Let $R$ be a system for which (LB) is already a theorem. Examples are resolution (Haken
1985; Urquhart 1987), polynomial calculus (Razborov 1998; Ben-Sasson--Impagliazzo 1999),
cutting planes (Pudl\'ak 1997), and bounded-depth Frege (Ajtai; Pitassi--Beame--Impagliazzo;
Kraj\'{\i}\v{c}ek--Pudl\'ak--Woods). For such $R$, (BR) holds if and only if
$\mathrm{P}\ne\mathrm{NP}$.
\end{theorem}
\begin{proof}
(LB) makes the consequent of (BR) false. So (BR) is equivalent to
$\neg(\mathrm{SAT}\in\mathrm{P})$.
\end{proof}
\noindent\emph{Consequence.} For every system whose hardness we can already prove, the
``weakest sufficient bridge'' is the full problem. In the vocabulary of the world ledger
(\texttt{research/agents/claude/world-ledger/}), it is an equivalent of $\mathrm{P}\ne\mathrm{NP}$.
It is not progress. A fixed multisign hypergraph calculus that is resolution-like,
algebraic or cutting-planes-like falls under Theorem~\ref{thm:reform}. Any proof of (BR)
for it would need a separate argument that is at least as strong as $\mathrm{P}\ne\mathrm{NP}$.

\begin{definition}[Correctness tautologies]
For a deterministic polynomial-time machine $A$ running in time $n^c$, let
$\mathrm{Corr}_{A,n}(\vec f,\vec x,\vec z)$ be the propositional formula of size
$\poly(n)$ stating the following. If $\vec z$ is the valid computation tableau of $A$ on
the length-$n$ input $\vec f$, and the tableau ends in REJECT, then the assignment
$\vec x$ falsifies the CNF encoded by $\vec f$. If $A$ decides SAT, every
$\mathrm{Corr}_{A,n}$ is a tautology.
\end{definition}

\begin{theorem}[PROVED/EXACT; standard reflection argument, cf.\ Cook 1975, Kraj\'{\i}\v{c}ek--Pudl\'ak 1989]\label{thm:corr}
Let $R$ p-simulate Extended Frege, and let $\Wb_R$ be measured on CNF refutations, where
tautologies are handled through the Tseitin encoding of their negations. Let $A$ be a
correct polynomial-time SAT decider. Then
\[
\Wb_R(N)\le\poly(N)\iff \mathrm{Corr}_{A,n}\text{ has $R$-proofs of size }\poly(n).
\]
Consequently,
\[
(\mathrm{BR})\iff\Big[\text{every correct poly-time SAT decider }A\text{ has }\poly\text{-size }R\text{-proofs of }\mathrm{Corr}_{A,n}\Big].
\]
\end{theorem}
\begin{proof}
($\Rightarrow$) $\mathrm{Corr}_{A,n}$ is a tautology of size $\poly(n)$. Polynomial
boundedness gives a short refutation of the Tseitin encoding of its negation.
($\Leftarrow$) Take an unsatisfiable $F$ of length $n$. Substitute the constants
$\langle F\rangle$ for $\vec f$ in the given proof. EF is closed under substitution with
polynomial overhead. Then derive the tableau $\vec z$ by evaluating the $\poly(n)$-size
circuit of $A$ gate by gate. Each gate needs $O(1)$ EF lines, and the extension
variables name the gates. Modus ponens gives $\neg F(\vec x)$. Combining this with the
clauses of $F$ yields $\bot$ in $O(|F|)$ further lines. For the second equivalence: if
SAT$\in$P via some $A$ and (BR) holds, then $R$ is p-bounded, so every correct $A'$
qualifies by ($\Rightarrow$). Conversely, if every correct $A$ qualifies and some correct
$A$ exists, ($\Leftarrow$) gives (BR).
\end{proof}

\begin{corollary}[CONDITIONAL; Cook--Kraj\'{\i}\v{c}ek 2007]
If $\mathrm{SAT}\in\mathrm{P}$ is \emph{provable in $\mathrm{PV}$} (equivalently, in
$S^1_2$), then (BR) holds for $R=\mathrm{EF}$. Cook's translation turns PV proofs of
$\forall$-statements into polynomial-size EF proofs of their propositional instances.
Hence (BR) for EF can fail only in the world where $\mathrm{P}=\mathrm{NP}$ holds but no
correct SAT algorithm is $\mathrm{PV}$-provably correct.
\end{corollary}

\paragraph{The exact open arrow.}
With $R\succeq\mathrm{EF}$, the directive needs \textbf{two} results. Both are OPEN:
\begin{enumerate}
\item[(LB$_{\mathrm{EF}}$)] Superpolynomial Extended Frege lower bounds. This is the
central open problem of proof complexity. It alone already implies
$\mathrm{PV}\nvdash\mathrm{P}=\mathrm{NP}$, but not $\mathrm{P}\ne\mathrm{NP}$.
\item[(CPA)] The correctness-provability arrow of Theorem~\ref{thm:corr}. Any correct
fast SAT algorithm must be correct ``for an EF-short reason''.
\end{enumerate}
The run of $A$ itself is \emph{cheap} to encode. Tableau evaluation is polynomial. All the
difficulty lives in the soundness lemma ``REJECT $\Rightarrow$ UNSAT''. This lemma is a
universal statement about $A$. A trace of the computation cannot supply it. So a
simulation that charges only the steps $A$ executes cannot prove (BR). It must charge the
\emph{proof of correctness}. This identifies precisely where earlier corpus transfer
attempts lose: R\_mix341, the FLMV machine, and the TMR cost-sound abstraction obligations.

\section{Hostile controls: pointwise transfer and strategy switching}
\begin{proposition}[KNOWN THEOREM + COMPUTATIONAL]\label{prop:tseitin}
Pointwise transfer $T_A(F)\ge\cost_R^{\min}(F)^{\Omega(1)}$ fails for $R=$ resolution,
and for every multisign reduction calculus that resolution p-simulates. Take Tseitin
formulas on bounded-degree expanders. Resolution needs size $2^{\Omega(n)}$ (Urquhart 1987;
width $\Omega(n)$ and size $\ge\exp((w-3)^2/\Omega(m))$ by Ben-Sasson--Wigderson 2001).
Gaussian elimination over GF(2) refutes them in $O(n^3)$ bit operations.
\end{proposition}
Table~\ref{tab:ts} reports exact minimum refutation widths from width-bounded saturation.
At widths 3 and 4 the closure saturates without $\bot$, so the bound ``$\ge5$'' is exact.
At width 5 the run stopped at the 25-second budget, so that row is ``undetermined''. The
same table gives the exact GF(2) elimination bit-operation count. Script:
\texttt{research/agents/claude/wcdc/tseitin\_width\_vs\_gf2.py}. Raw data:
\texttt{tseitin\_results.json}. Chunking: one graph per invocation, with a 25\,s budget
per width level.
\begin{table}[h]\centering
\begin{tabular}{lrrrll}\toprule
Graph & vars & clauses & GF(2) bit-ops & min width & width-$w$ closure sizes\\\midrule
$K_4$ & 6 & 16 & 35 & $=4$ & 16, 473\\
$K_{3,3}$ & 9 & 24 & 110 & $=5$ & 24, 384, 6883\\
$Q_3$ & 12 & 32 & 260 & $\ge5$ & 32, 512, ($>30348$ at $w{=}5$, budget)\\
Petersen & 15 & 40 & 336 & $\ge5$ & 40, 640, ($>14240$ at $w{=}5$, budget)\\\bottomrule
\end{tabular}
\caption{COMPUTATIONAL. Width grows above the initial width 3, while elimination stays
cubic. At these sizes this is an illustration only. The asymptotic separation is the
cited theorem.}\label{tab:ts}
\end{table}

\paragraph{Adaptive strategy switching.}
A union $R_1\cup R_2$, where a proof is a proof in either system, is again a
Cook--Reckhow system with $\cost^{\min}_{R_1\cup R_2}=\min$. Hardness must survive the
union. The known lower-bound families do not:
\begin{center}
\begin{tabular}{lll}\toprule
Family & Hard for & Easy for\\\midrule
Tseitin (expanders) & Res, PC/$\mathbb F_p$ ($p\ne2$) & PC/$\mathbb F_2$, Res($\oplus$), GF(2) elimination; CP quasi-poly (Dadush--Tiwari 2020)\\
PHP$^{n+1}_n$ & Res, bounded-depth Frege, PC & Cutting planes, Frege\\
Random 3-CNF, $m=\Delta n$ & Res, PC (Ben-Sasson--Impagliazzo) & no known short proofs (CANDIDATE)\\\bottomrule
\end{tabular}
\end{center}
Random 3-CNF is the only standard \textbf{CANDIDATE} family that survives the combined
controls. For bounded-depth Frege only an $\Omega(n^{1+\varepsilon_k})$ \emph{superlinear}
bound is known (Gryaznov--Talebanfard, arXiv 2403.02275). DAG-like Res($\oplus$) has no
superpolynomial lower bound. Only the regular fragment (Efremenko--Garl\'{\i}k--Itsykson,
SICOMP 2025) and depth-restricted fragments (Alekseev--Itsykson; Bhattacharya--Chattopadhyay;
Byramji--Impagliazzo) have one. Even if random 3-CNF were EF-hard, Theorem~\ref{thm:corr}
shows that the bridge would still be missing.

\section{Barrier and transfer audit}
\begin{itemize}
\item \textbf{Relativization / algebrization.} Theorem~\ref{thm:reform} is
oracle-independent. (CPA) is a statement about provability. There are oracle worlds in
which $\mathrm{P}=\mathrm{NP}$, and nothing in (CPA) is known to relativize favorably.
FAILED TRANSFER for any oracle-uniform argument.
\item \textbf{Natural proofs.} An EF lower bound by a ``natural'' hardness predicate
faces the Razborov--Rudich-type obstruction for strong systems. Feasible interpolation
fails for EF under standard cryptographic assumptions (Kraj\'{\i}\v{c}ek--Pudl\'ak 1998).
\item \textbf{PQ ledger item 6 (special soundness).} Theorem~\ref{thm:corr} is the exact
proof-complexity analogue. The ``binding'' of the decider's REJECT bit to unsatisfiability
must be \emph{proved}; it cannot be assumed. EXACT OBJECT TRANSFER of the structural
lesson. It brings no new theorem.
\item \textbf{Electronics / verification analogy.} (CPA) is the claim ``every correct
circuit family has a polynomial-size equivalence certificate in a fixed proof format''.
Hardware verification knows this certificate gap: miters that are correct but have no
short resolution proof. USEFUL ANALOGY only. It needs a transfer to EF.
\end{itemize}

\section{Corpus reviewed; duplication and contradiction audit}
\textbf{CORPUS REVIEWED.} Repository tree and recent \texttt{git log} (through TMR-3/043 and
GROKPROOF 12:20 UTC); \texttt{AGENTS.md}; the TMR replacement directive
\texttt{research/coordination/2026-11-13\_1249\_TM\_RESOURCE\_AGENT\_REPLACEMENT.md};
the codex reports \texttt{research/agents/codex/2026-11-11\_*} (SAT polynomial assumption
audit, transfer recovery); the world ledger; repository search for
$\Wb_R$, ``worst-case derivation'' and ``weakest sufficient''.\\
\textbf{DUPLICATION / CONTRADICTION AUDIT.} The search found no prior artifact defining
$\Wb_R(N)$ under this directive. LSPV-1/350 (MUS worst-case quotient) and the TMR-1 rule
that rejects per-length minimization over \emph{algorithms} are consistent. Here we
minimize over \emph{derivations} of one fixed $R$, which has the correct quantifiers.
There is no contradiction with any root tracker claim.\\
\textbf{WHAT CHANGED RELATIVE TO ROOT MAIN.} We added the directive record and the
classification Theorems~\ref{thm:reform}--\ref{thm:corr}. Root \texttt{main.tex} is not
edited: this routine owns no marked region there, and LSPV-5/TMR-5 is synthesis arbiter.\\
\textbf{ADOPTED / REJECTED / RETIRED.} Adopted: the fixed-$R$ worst-case formulation.
Rejected as a weaker step: (BR) for any $R$ with a known lower bound (it is a
reformulation). Retired for this route: pointwise algorithm-trace simulation into
resolution-type multisign calculi (Proposition~\ref{prop:tseitin}).\\
\textbf{STALE OR OVERSTATED CLAIMS FOUND.} The directive's phrase ``weakest sufficient
bridge'' suggests the bridge is easier than $\mathrm{P}\ne\mathrm{NP}$. By
Theorem~\ref{thm:reform} it is not easier for any currently hard $R$, and by
Theorem~\ref{thm:corr} it is a specific open provability statement for strong $R$.

\section{Status ledger and next target}
\begin{center}
\begin{tabular}{ll}\toprule
Arrow & Status\\\midrule
$\Wb_R$ = Cook--Reckhow length; free optimal sequence is legitimate & KNOWN THEOREM\\
(BR) $\Leftrightarrow$ P$\ne$NP for $R\in\{$Res, PC, CP, AC$^0$-Frege$\}$ & PROVED/EXACT\\
(BR) $\Leftrightarrow$ (CPA) for $R\succeq$ EF & PROVED/EXACT\\
PV $\vdash$ SAT$\in$P $\Rightarrow$ (BR) for EF & CONDITIONAL (known)\\
Pointwise transfer into Res-type calculi & FALSIFIED/RETIRED\\
(LB$_{\mathrm{EF}}$) & OPEN\\
(CPA) & OPEN\\
Random 3-CNF hard for combined systems & CANDIDATE\\\midrule
P versus NP & \textbf{OPEN}\\\bottomrule
\end{tabular}
\end{center}
\textbf{Next highest-value target.} Find either a \emph{counterexample shape} or a
\emph{partial proof} for (CPA). Ask whether there is a class of correct polynomial-time
SAT procedures broader than PV-definable-and-provable for which EF-short correctness
proofs are forced: for example, algorithms whose correctness follows from a
bounded-arithmetic-provable invariant on the multisign hypergraph state. The red-team
counterpart is a correct polynomial-time decider on a promise family (e.g.\ Tseitin via
elimination) whose correctness tautologies are provably hard for a fixed
$R\not\succeq\mathrm{EF}$. Res and $\mathrm{Corr}$ for Gaussian elimination is the first
test.

\paragraph{Build note.} No TeX engine is installed in this container, so this file was
checked structurally (balanced environments) and was not compiled.

\section*{References}
\small
S.~Cook, R.~Reckhow, The relative efficiency of propositional proof systems, JSL 44 (1979).
S.~Cook, Feasibly constructive proofs and the propositional calculus, STOC 1975.
S.~Cook, J.~Kraj\'{\i}\v{c}ek, Consequences of the provability of NP$\subseteq$P/poly, JSL 72 (2007).
J.~Kraj\'{\i}\v{c}ek, P.~Pudl\'ak, Propositional proof systems, the consistency of first order theories and the complexity of computations, JSL 54 (1989).
A.~Haken, The intractability of resolution, TCS 39 (1985).
A.~Urquhart, Hard examples for resolution, JACM 34 (1987).
E.~Ben-Sasson, A.~Wigderson, Short proofs are narrow, JACM 48 (2001).
E.~Ben-Sasson, R.~Impagliazzo, Random CNFs are hard for the polynomial calculus, FOCS 1999.
P.~Pudl\'ak, Lower bounds for resolution and cutting plane proofs and monotone computations, JSL 62 (1997).
S.~Gryaznov, N.~Talebanfard, Bounded-depth Frege lower bounds for random 3-CNFs via deterministic restrictions, arXiv:2403.02275.
K.~Efremenko, M.~Garl\'{\i}k, D.~Itsykson, Lower bounds for regular resolution over parities, SICOMP 2025.
D.~Dadush, S.~Tiwari, On the complexity of branching proofs, CCC 2020.

% END_OF_VERIFIED_REPORT::CLAUDE::WCDC-0001
\end{document}
